% !PRINTED 481
\weberpage{481}
%<<CONT
For the coefficients $a_1, a_2 \ldots a_m$, regarded as independent
variables, can also be represented as the symmetric fundamental functions of
the roots $\alpha, \alpha_1, \alpha_2 \ldots \alpha_{m-1}$. If then $g(t)$ is
an irreducible factor of $G(t)$, expressed rationally in terms of the $a$,
which vanishes for $t = \varrho$, then, on representing $a_i$ and $\varrho$
by the $\alpha$, one obtains in $g(\varrho) = 0$ an identical equation. In
this, however, the $\alpha$ may be permuted arbitrarily, whereby the $a$ do
not change, while $\varrho$ can pass into any of the roots
$\varrho, \varrho', \varrho'' \ldots$ of $G(t)$. Hence $g(t)$ vanishes for
all the roots of $G(t)$ and must consequently coincide with
$G(t)$\footnote{\'Evariste \gesperrt{Galois} fell in a duel in the
year 1832, scarcely 20 years old. The first indication of the theory that
today bears his name is found in a memoir ``Analyse d'un m\'emoire sur la
r\'esolution alg\'ebrique des \'equations'', which appeared in 1830 in the
``Bulletin des sciences math\'em.''\ of \gesperrt{F\'erussac}. More detailed
communications are contained in the famous letter to \gesperrt{Auguste Chevalier},
written on the eve of his death, which is published in the ``Revue
encyclop\'edique'' of September 1832. Only in the year 1846 did \gesperrt{Liouville}
have all the already published papers of \gesperrt{Galois}, together with some
investigations from his literary remains, among them the most important,
``M\'em. sur les conditions de r\'esolubilit\'e des \'equations par
radicaux'', printed in Vol.~XI of his Journal. A German edition by \gesperrt{Maser}
appeared in 1889 (Berlin, published by Springer). Interesting biographical
communications are found in the ``Revue encyclop\'edique'' of September
1832, which \gesperrt{Liouville} did not reprint along with them and which are also
lacking in the German edition. Of further writings that have contributed to
the understanding or the further development of the \gesperrt{Galois} theory, let there
still be mentioned: J. A. \gesperrt{Serret} (``Cours d'alg\`ebre sup\'erieure'',
II. edition, 1854, IV. 1879); \gesperrt{Betti} (Annali di scienze fisiche e matematiche
[1853]); C. \gesperrt{Jordan}, ``Trait\'e des substitutions'' (1870). E. \gesperrt{Netto},
``Substitutionentheorie'' (Leipzig 1882).}.

\begin{center}
\S.\,150.

\gesperrt{Transitive and intransitive groups.}
\end{center}

From the \gesperrt{Galois} group $P$ we can derive a very
simple criterion for whether the equation $F(x) = 0$ is reducible or
irreducible.

Let us assume that $F(x)$, of degree $m$, is reducible, and that $f(x)$ is a
factor of $F(x)$ in $\Omega$ of degree $n$, with $n < m$. Let the roots of
$f(x)$
%CONT>>
